\chapter{Renormalization Group}

\begin{remarkbox}
Renormalization introduced a scale \(\mu\). The renormalization group explains
how physics remains independent of this arbitrary scale while renormalized
parameters run with it. It is the language of scale dependence in quantum field
theory.
\end{remarkbox}

\section{Why a Group of Scale Transformations?}

In the previous chapter, renormalized parameters were defined at a scale
\(\mu\). A coupling written as \(g(\mu)\) is not a fixed number independent of
context. It depends on the scale at which it is measured or defined.

This does not mean that physical predictions depend on an arbitrary convention.
A physical observable \(\mathcal{O}\) may depend explicitly on \(\mu\) through
loop logarithms, and implicitly on \(\mu\) through the running parameters:
\[
    \mathcal{O}=\mathcal{O}(p_i,m(\mu),g(\mu),\mu).
\]
The total dependence on \(\mu\) must cancel:
\[
    \mu\frac{\dd}{\dd\mu}\mathcal{O}=0.
\]
The renormalization group is the formalism that organizes this cancellation.

The word ``group'' reflects the fact that changing from one scale to another can
be composed: running from \(\mu_1\) to \(\mu_2\), and then from \(\mu_2\) to
\(\mu_3\), is equivalent to running directly from \(\mu_1\) to \(\mu_3\).

\section{Running Couplings}

A running coupling satisfies an equation of the form
\[
    \mu\frac{\dd g}{\dd\mu}=\beta(g).
\]
The function \(\beta(g)\) is the beta function. It is determined by quantum loop
corrections and depends on the theory and renormalization scheme.

If \(\beta(g)>0\), the coupling grows as the scale increases. If
\(\beta(g)<0\), the coupling decreases as the scale increases. The sign and
zeros of the beta function therefore encode important physical information.

For QED with one charged fermion, the leading beta function is
\[
    \beta(e)=\frac{e^3}{12\pi^2}+\cdots.
\]
The positive sign means that the effective electric charge grows slowly at short
distances.

\section{Solving a Simple RG Equation}

Suppose the beta function has the form
\[
    \mu\frac{\dd g}{\dd\mu}=b g^3.
\]
Then
\[
    \frac{\dd g}{g^3}=b\frac{\dd\mu}{\mu}.
\]
Integrating between \(\mu_0\) and \(\mu\) gives
\[
    \frac{1}{g(\mu)^2}
    = \frac{1}{g(\mu_0)^2}-2b\ln\frac{\mu}{\mu_0}.
\]
This equation shows how logarithms are reorganized into scale-dependent
couplings.

If \(b>0\), the coupling grows with \(\mu\). If \(b<0\), the coupling becomes
weaker at high energies. The latter behavior is called asymptotic freedom.

\section{The Callan--Symanzik Equation}

Correlation functions also satisfy renormalization group equations. For an
\(n\)-point renormalized correlation function \(G^{(n)}\), the
Callan--Symanzik equation has the schematic form
\[
    \left[
      \mu\frac{\partial}{\partial\mu}
      +\beta(g)\frac{\partial}{\partial g}
      +m\gamma_m(g)\frac{\partial}{\partial m}
      +n\gamma_\phi(g)
    \right]G^{(n)}=0.
\]
Here \(\gamma_\phi\) is the anomalous dimension of the field, and
\(\gamma_m\) is the anomalous dimension of the mass.

The equation says that explicit scale dependence is compensated by the running
of couplings, masses, and field normalizations. It is a precise version of the
statement that physics does not depend on the arbitrary renormalization scale.

\section{Anomalous Dimensions}

Classically, a field has an engineering dimension determined by dimensional
analysis. For a scalar field in four dimensions,
\[
    [\phi]_{\mathrm{classical}}=1.
\]
Quantum corrections can modify scaling behavior. The correction is called an
anomalous dimension.

If \(Z_\phi\) is the field-strength renormalization, the anomalous dimension is
schematically
\[
    \gamma_\phi(g)=\frac{1}{2}\mu\frac{\dd}{\dd\mu}\ln Z_\phi.
\]
The full scaling dimension is then
\[
    \Delta_\phi=\Delta_{\mathrm{classical}}+\gamma_\phi.
\]
Anomalous dimensions are especially important near fixed points, where they give
universal critical exponents.

\section{Fixed Points}

A fixed point is a value \(g_*\) for which
\[
    \beta(g_*)=0.
\]
At a fixed point, the coupling no longer runs. The theory may become scale
invariant, and often conformally invariant under additional assumptions.

The Gaussian fixed point is
\[
    g_*=0.
\]
It describes a free theory. Interacting fixed points, if they exist, describe
scale-invariant interacting quantum field theories.

The behavior near a fixed point is determined by linearizing the beta function:
\[
    \beta(g)\approx \beta'(g_*)(g-g_*).
\]
If perturbations grow as one moves to low energies, they are relevant. If they
shrink, they are irrelevant. If they remain of comparable size, they are
marginal.

\section{Relevant, Irrelevant, and Marginal Operators}

Operators are classified by how their couplings scale. In four dimensions, an
operator \(\mathcal{O}\) of dimension \(\Delta\) appears as
\[
    \mathcal{L}\supset c\,\mathcal{O}.
\]
The coupling has dimension
\[
    [c]=4-\Delta.
\]
If \(\Delta<4\), the operator is relevant. If \(\Delta=4\), it is marginal. If
\(\Delta>4\), it is irrelevant.

Relevant operators become more important in the infrared. Irrelevant operators
are suppressed at low energies by powers of a high scale. Marginal operators
require quantum analysis: their beta functions decide whether they are marginally
relevant, marginally irrelevant, or exactly marginal.

This classification is the bridge between renormalization group and effective
field theory.

\section{Dimensional Transmutation}

A theory may begin with a dimensionless coupling but generate a physical mass
scale through running. This phenomenon is called dimensional transmutation.

Suppose
\[
    \mu\frac{\dd g}{\dd\mu}=-b g^3,
    \qquad b>0.
\]
Then the running coupling can be written in terms of a scale \(\Lambda\):
\[
    g(\mu)^2\sim \frac{1}{2b\ln(\mu/\Lambda)}.
\]
The dimensionless coupling has been exchanged for a dimensionful scale
\(\Lambda\). This idea is central in non-Abelian gauge theories and QCD, where a
strong-coupling scale emerges from a classically scale-invariant gauge coupling.

\section{Asymptotic Freedom}

A theory is asymptotically free if its coupling becomes small at high energies:
\[
    g(\mu)\longrightarrow0
    \qquad \text{as} \qquad \mu\to\infty.
\]
This occurs when the beta function is negative at small coupling:
\[
    \beta(g)=-b g^3+\cdots,
    \qquad b>0.
\]
Non-Abelian gauge theories can have this property. It is the reason quarks and
gluons behave approximately as weakly interacting particles at very short
distances, even though they are confined at long distances.

QED is not asymptotically free. Its coupling grows slowly at short distances.
This contrast between Abelian and non-Abelian gauge theories is one of the most
important lessons of the renormalization group.

\section{Landau Poles}

If a coupling grows without bound at a finite scale in perturbation theory, the
scale is called a Landau pole. In a simple running equation with positive
\(b\),
\[
    \frac{1}{g(\mu)^2}
    = \frac{1}{g(\mu_0)^2}-2b\ln\frac{\mu}{\mu_0},
\]
the denominator can vanish at finite \(\mu\). Perturbation theory then breaks
down.

A Landau pole may signal that the theory cannot be extended to arbitrarily high
energies without new physics, or that perturbation theory is no longer reliable.
In effective field theory, this is not necessarily a problem; it simply indicates
the scale at which the description must be replaced or completed.

\section{RG Improvement of Perturbation Theory}

Perturbative calculations often contain large logarithms such as
\[
    \lambda\ln\frac{E}{\mu}.
\]
If \(E\) is very different from \(\mu\), these logarithms can spoil the
convergence of perturbation theory even when \(\lambda\) is small.

The renormalization group improves perturbation theory by choosing \(\mu\) near
the physical scale and using running couplings. This resums towers of logarithms
into scale-dependent parameters.

The practical rule is:
\[
    \text{choose } \mu \text{ of order the characteristic energy scale}.
\]
Then the explicit logarithms are small, and the running coupling carries the
large-scale dependence.

\section{Physical Meaning of Running}

Running couplings do not mean that charge or interaction strength changes with
time. They mean that different experiments probe the theory at different length
or momentum scales. A coupling is an effective parameter summarizing the response
of the quantum vacuum at a chosen scale.

At long distances, high-energy fluctuations have already been integrated into
the parameters. At short distances, more fluctuations are resolved explicitly.
The renormalization group describes how the description changes as the resolution
changes.

This is the conceptual heart of modern QFT: the same physics can be described at
many scales, but the parameters and relevant degrees of freedom may change.

\begin{summarybox}
The renormalization group organizes the dependence of renormalized quantities on
the scale \(\mu\). Couplings run according to beta functions, fields and masses
acquire anomalous dimensions, and physical observables remain independent of the
arbitrary renormalization scale. Fixed points, relevant and irrelevant operators,
dimensional transmutation, asymptotic freedom, and RG improvement are all
consequences of scale-dependent quantum fluctuations.
\end{summarybox}

\section{Chapter Checklist}

After reading this chapter, one should be able to:
\begin{itemize}
    \item explain why renormalized parameters depend on \(\mu\);
    \item define a beta function;
    \item solve a simple one-coupling RG equation;
    \item state the meaning of the Callan--Symanzik equation;
    \item explain anomalous dimensions;
    \item define fixed points;
    \item classify relevant, irrelevant, and marginal operators;
    \item explain dimensional transmutation;
    \item distinguish asymptotic freedom from Landau-pole behavior;
    \item explain how RG improvement resums large logarithms;
    \item describe the physical meaning of running couplings.
\end{itemize}
